\section{Related literature}

Our closest reference is \citet{greenwood2025}, who study the declining S\&P 500 index effect and examine whether changes in demand, the composition of index events and trading conditions account for it. We apply these empirical questions to OMXSB reconstitutions, adapting the benchmark, sample period and holdings measures to Swedish data. Our timing extension examines how the historical comparison changes when the announcement window has a fixed duration and uses the same stocks as the full event window. We also compare holdings measurements across two reporting sources and follow cumulative returns on a fixed event sample after implementation. These exercises connect the return comparison to its institutional timing and measured trading counterpart.

The interpretation of index-event returns builds on \citet{shleifer1986} and \citet{harris1986}. The former connects index additions to downward-sloping demand for individual stocks, while the latter examines price pressure and subsequent reversal. These studies motivate looking beyond the announcement return to the timing of trading and the subsequent price path. A positive cumulative return at a later date is relevant to that comparison, but does not by itself distinguish a lasting demand response from information or selection associated with index membership.

\citet{chang2015} study the price effects of indexing through a regression discontinuity design. Their setting illustrates the additional identifying structure needed to interpret an index comparison causally. OMXSB's selection criteria differ, and the present study does not reconstruct an equivalent assignment boundary. Our composition adjustments and comparison securities therefore describe conditional return differences rather than treatment effects.

Finally, \citet{chinco2024} show why the investor population relevant to passive ownership can be broader than explicitly labeled index funds. This distinction matters for measuring the demand associated with OMXSB changes, where benchmark variants and benchmark-linked mixed funds can hold the same securities. We consequently distinguish observed changes in the covered funds' positions from total passive demand. The combined contribution is evidence on the return pattern and its trading counterpart in an additional institutional setting, with a direct measurement comparison for reported holdings.
